% !TEX root = ../main.tex

% 分析与讨论	
\begin{table}
	\renewcommand\arraystretch{1.7}
	\begin{tabularx}{\textwidth}{|X|X|X|X|}
		\hline
		Major：& Physics &Grade：& 2022\\
		\hline
		Name： & 杨舒云 & Student number：& 22344020\\
		\hline
		Date：& 2024/9/27 & Score： &\\
		\hline
	\end{tabularx}
\end{table}
\section{APL1-3 Quantum Key Distribution \\ Analysis \& Discussion}


%---------------------------------------------------------------------
% 数据处理
\subsection{Data Processing}

% 分析
\subsubsection{Analysis}
\begin{enumerate}
	\item 计算单光子探测器的探测效率
	
	根据\cref{tab:tab2-1}和\cref{tab:tab2-2}，我们做如下计算：
	
	404nm 单光子的能量为：
	$$ E = h \nu = \frac{h c}{\lambda} = 6.626 \times 10^{-34} \times \frac{3 \times 10^{8}}{404 \times 10^{-9}} \mathrm{J} = 4.92 \times 10^{-19} \mathrm{J} $$
	
	404nm 激光器的发光重复频率为 10MHz，当到达$\mathrm{SPD_M}$的每个脉冲平均光子数为 0.1 个光子时，$\mathrm{SPD_M}$ 所探测到的功率为：
	\begin{align*}
		P = \mu \cdot f \cdot E = 0.1 \times 10 \times 10^{6} \times 4.92 \times 10^{-19} \mathrm{W} = 4.92 \times 10^{-13} \mathrm{W}	\\
		= 10 \lg \frac{4.92 \times 10^{-13} \mathrm{W}}{1 \times 10^{-3} \mathrm{W}} \mathrm{dBm} = -93.08 \mathrm{dBm}
	\end{align*}
	其中 $\mu$ 为平均光子数脉冲，$f$ 为光触发频率，$P$ 即为光功率。
	
	使用手持光功率计测量 $\mathrm{LD_M}$ 的光功率为 $P_0 = 10.7 \mathrm{dBm}$，所以为制备每个脉冲平均光子数为 0.1 个光子，在路径中需要加的总衰减为：
	$$
	A = 103.78\text{dB}
	$$
	
	固定衰减（注意到光纤耦合器 $\mathrm{OC_M}$ 的光纤耦合效率所提供的衰减 48\% ）：
	$$
	A_1 = 93.7\text{dB} + 10\lg(\frac{1}{0.48}) = 96.89\text{dB}
	$$
	
	由此得可调衰减：
	$$
	A_2 = A - A_1 = 6.89\text{dB}
	$$
	
	于是，由$N = f\mu\eta$，推得探测器效率：
	$$
	\eta = \frac{N}{f\mu} = \frac{353000}{10\times10^6\times0.1} = 35.3\%
	$$
	
	\item 计算单光子标定
	
	理论计算当公共信道平均光子数为 0.2 光子/脉冲时，经过接收端光路衰减到达单光子探测器的探测计数 N，（以 M 路光路为例计算的）计算过程如下：
	
	当公共信道处光功率为：
	\begin{align*}
		P_{pc} = 9.84 \times 10^{-13} \mathrm{W} = 10 \lg \frac{9.84 \times 10^{-13} \mathrm{W}}{1 \times 10^{-3} \mathrm{W}} \mathrm{dBm} = -90.07 \mathrm{dBm}
	\end{align*}
		
	公共信道处的光子到达 $\mathrm{OC_M}$ 后，此时又经过了4个衰减，那么最后的功率为：
	\begin{align*}
		P = -96.66 \mathrm{dBm} = 10^{\frac{-96.66}{10}} \mathrm{mW} = 2.16 \times 10^{-10} \mathrm{mW} = 2.16 \times 10^{-13} \mathrm{W}
	\end{align*}	
		
	到达 $\mathrm{OC_M}$ 的每脉冲平均光子数为：
	$$ \mu = \frac{P}{f \cdot E} = \frac{2.16 \times 10^{-13}}{10 \times 10^{6} \times 4.92 \times 10^{-19}} = 0.044 $$
		
	最后得单光子探测器的探测计数 N 为：
	$$ N = f \cdot \mu \cdot \eta = 10 \times 10^{6} \times 0.044 \times 35.3 \% = 155320 $$
	
	\item 关于误码率的计算
	\begin{itemize}
		\item 先来看看一般情况：
		
		要利用十六进制的符号误码率近似计算二进制的比特误码率，我们可以从概率论的角度进行详细推导。首先，定义比特误码率 \( P_b \) 为单个比特发生错误的概率，符号误码率 \( P_s \) 为一个符号发生错误的概率。在十六进制系统中，一个符号由4个比特组成，因此如果比特误码是独立的，那么一个符号正确接收的概率就是组成该符号的所有比特都正确接收的概率，即：
		
		\[
		P_{\text{正确符号}} = (1 - P_b)^4
		\]
		
		因此，符号误码率可以表示为符号接收错误的概率，即至少有一个比特出错的概率：
		
		\[
		P_s = 1 - P_{\text{正确符号}} = 1 - (1 - P_b)^4
		\]
		
		接下来，我们对 \( (1 - P_b)^4 \) 进行展开，并利用近似方法进行计算。利用二项式定理展开，得到：
		
		\[
		(1 - P_b)^4 = 1 - 4P_b + 6P_b^2 - 4P_b^3 + P_b^4
		\]
		
		将展开式代入符号误码率的表达式，得到：
		
		\[
		P_s = 1 - \left(1 - 4P_b + 6P_b^2 - 4P_b^3 + P_b^4\right) = 4P_b - 6P_b^2 + 4P_b^3 - P_b^4
		\]
		
		由于 \( P_b \) 通常是很小的值，可以忽略高次幂项（如 \( P_b^2 \)、\( P_b^3 \)、\( P_b^4 \)），因此符号误码率近似为：
		
		\[
		P_s \approx 4P_b
		\]
		
		由此，我们可以得到比特误码率与符号误码率之间的近似关系：
		
		\[
		P_b \approx \frac{P_s}{4}
		\]
		
		如果需要更精确的结果，可以保留更多的项。例如，保留二次项：
		
		\[
		P_s \approx 4P_b - 6P_b^2
		\]
		
		此时，可以解二次方程 \( 6P_b^2 - 4P_b + P_s = 0 \) 来求得更准确的 \( P_b \) 值。但对于较小的 \( P_b \) 值，一次近似已经足够准确。
		
		\item 而对于本实验中的十六进制文件，我们具体观察（如\cref{fig:example_of_hex}所示）便可知道，一个符号由两个十六进制数组成，也就是说只要对应的8个比特中任一个出错，那么整个16进制符号就出错了，因此要利用由两个十六进制数（总共8个比特）组成的符号误码率来近似计算二进制的比特误码率，我们同样可以从概率论的角度进行新的推导。
		
		\begin{figure}[htbp]
			\centering
			\includegraphics[width=0.5\textwidth]{APL1_3_example_of_hex.png}
			\caption{十六进制文件格式样例}
			\label{fig:example_of_hex}
		\end{figure}
		
		首先，定义比特误码率 \( P_b \) 为单个比特发生错误的概率，符号误码率 \( P_s \) 为一个符号发生错误的概率。在这种情况下，一个符号由8个比特组成，因此，如果比特误码是独立的，一个符号正确接收的概率就是组成该符号的所有比特都正确接收的概率，即：
		
		\[
		P_{\text{正确符号}} = (1 - P_b)^8
		\]
		
		符号误码率则表示为符号接收错误的概率，即至少有一个比特出错的概率：
		
		\[
		P_s = 1 - P_{\text{正确符号}} = 1 - (1 - P_b)^8
		\]
		
		接下来，对 \( (1 - P_b)^8 \) 进行展开。利用二项式定理展开，得到：
		
		\[
		(1 - P_b)^8 = 1 - 8P_b + 28P_b^2 - 56P_b^3 + 70P_b^4 - 56P_b^5 + 28P_b^6 - 8P_b^7 + P_b^8
		\]
		
		将展开式代入符号误码率的表达式，得到：
		
		\[
		\begin{aligned}
			P_s &= 1 - \left(1 - 8P_b + 28P_b^2 - 56P_b^3 + 70P_b^4 - 56P_b^5 + 28P_b^6 - 8P_b^7 + P_b^8\right) \\
			&= 8P_b - 28P_b^2 + 56P_b^3 - 70P_b^4 + 56P_b^5 - 28P_b^6 + 8P_b^7 - P_b^8
		\end{aligned}
		\]
		
		由于 \( P_b \) 通常是很小的值，可以忽略高次幂项（如 \( P_b^3 \) 及以上的项），因此符号误码率近似为：
		
		\[
		P_s \approx 8P_b - 28P_b^2
		\]
		
		为了得到比特误码率与符号误码率之间的关系，我们可以解这个近似的二次方程。将 \( P_s \) 表示为关于 \( P_b \) 的方程：
		
		\[
		28P_b^2 - 8P_b + P_s = 0
		\]
		
		解这个二次方程，得到：
		
		\[
		P_b = \frac{8 \pm \sqrt{64 - 112P_s}}{56}
		\]
		
		由于 \( P_b \) 是一个概率，且 \( P_b \) 很小，我们选择较小的正根：
		
		\[
		P_b \approx \frac{8 - \sqrt{64 - 112P_s}}{56}
		\]
		
		当 \( P_s \) 很小时，\( \sqrt{64 - 112P_s} \approx 8 - \frac{112P_s}{16} = 8 - 7P_s \)，因此：
		
		\[
		P_b \approx \frac{8 - (8 - 7P_s)}{56} = \frac{7P_s}{56} = \frac{P_s}{8}
		\]
		
		因此，在 \( P_s \) 很小的情况下，比特误码率近似为符号误码率的八分之一：
		
		\[
		P_b \approx \frac{P_s}{8}
		\]
		
		如果需要更精确的结果，可以使用原始的二次方程解。例如，对于较大的 \( P_s \) 值，直接使用二次方程求解 \( P_b \) 会更加准确。
		
		\item 结合具体实验数据来看：
		
		由\cref{tab:tab4_hex}得十六进制误码率，利用上述公式计算二进制误码率并与其它测量值相比对结果如\cref{tab:tab_BRE_compared}所示。
		
		\begin{table}[htbp]
			\centering
			\begin{tabular}{|l|c|}
				\hline
				项目 & 数据 \\
				\hline
				16进制误码率 & 9.85\% \\
				\highlight{近似二进制误码率} & \highlight{1.23\%} \\
				测量值（由QKD软件统计） & 1.20\% \\
				测量值（由Python脚本统计） & 1.29\% \\
				\hline
			\end{tabular}
			\caption{二进制误码率得计算比对}
			\label{tab:tab_BRE_compared}
		\end{table}
		
	\end{itemize}
\end{enumerate}

% 讨论
\subsubsection{Discussion}
\begin{enumerate}
	\item 关于“实验1：光的偏振的测量和控制”的讨论
	\begin{itemize}
		\item \textbf{假设激光器发出的光是线偏振光，求激光器发出的线偏振光与水平方向的夹角是多少？}
		
		根据\cref{tab:tab1}中测量得到的激光通过 PBS1 后得反射和投射功率，注意到对于一块偏振分光棱镜来说，它的透射量将是水平偏振的，它的反射量将是垂直偏振的，因此若是设激光的偏振方向与水平方向夹角为$\theta$，那么有：
		$$
		\frac{P_t}{P_r} = \frac{\cos^2\theta}{\sin^2\theta} = \cot^2\theta = 3.92 \Rightarrow \theta = 26.81^{\circ}
		$$
		此即偏振方向与水平方向夹角。
		
		\item \textbf{通过测量推导该激光器是否完美的线偏振？}
		
		根据测量数据\cref{tab:tab1}和如下所示得详细的理论计算分析，我们可以推导出该激光器发出的光是\textbf{高度线偏振}的，但\textbf{并非完美}的线偏振光。
		\begin{itemize}
			\item 首先，在没有半波片（HWP1）的情况下，反射功率 \(P_r\) 为 9.12 dBm（8.17 mW），透射功率 \(P_t\) 为 15.05 dBm（31.99 mW），总功率 \(P_{\text{total}}\) 为 \(P_r + P_t = 8.17\,\text{mW} + 31.99\,\text{mW} = 40.16\,\text{mW}\)。
			
			当调整半波片（HWP1）以获得最大透射功率时，最大透射功率 \(P_{t,\text{max}}\) 为 15.8 dBm（38.02 mW），最小反射功率 \(P_{r,\text{min}}\) 为 -1.66 dBm（0.68 mW），总功率为 \(38.02\,\text{mW} + 0.68\,\text{mW} = 38.70\,\text{mW}\)。
			
			当调整半波片以获得最大反射功率时，最大反射功率 \(P_{r,\text{max}}\) 为 15.88 dBm（38.73 mW），最小透射功率 \(P_{t,\text{min}}\) 为 -8.42 dBm（0.14 mW），总功率为 \(38.73\,\text{mW} + 0.14\,\text{mW} = 38.87\,\text{mW}\)。
			
			\item 理论上，偏振分束器（PBS1）会将入射光分成两个正交的偏振分量，透射功率 \(P_t\) 与 \(\cos^2\theta\) 成正比，反射功率 \(P_r\) 与 \(\sin^2\theta\) 成正比，其中 \(\theta\) 是入射光的偏振方向与 PBS1 透射轴之间的角度。半波片（HWP1）可以使入射光的偏振方向旋转 \(2\theta\)，从而控制进入 PBS1 的偏振态。
			
			\item 在分析过程中，首先计算无半波片时的偏振角度。根据测量数据，计算得到 \(\cos^2\theta = \frac{P_t}{P_{\text{total}}} = \frac{31.99}{40.16} \approx 0.7964\)，\(\sin^2\theta = \frac{P_r}{P_{\text{total}}} = \frac{8.17}{40.16} \approx 0.2036\)。因此，\(\tan^2\theta = \frac{\sin^2\theta}{\cos^2\theta} \approx \frac{0.2036}{0.7964} \approx 0.2557\)，进而得到 \(\theta = \arctan(\sqrt{0.2557}) \approx 26.6^\circ\)。这表明激光器的偏振方向与 PBS1 的透射轴之间的角度约为 \(26.6^\circ\)。
			
			接下来，计算消光比。透射功率的消光比为 \(\text{消光比}_{t} = \frac{P_{t,\text{max}}}{P_{t,\text{min}}} = \frac{38.02}{0.14} \approx 271:1\)，反射功率的消光比为 \(\text{消光比}_{r} = \frac{P_{r,\text{max}}}{P_{r,\text{min}}} = \frac{38.73}{0.68} \approx 57:1\)。这些有限的消光比表明激光器的偏振光具有高偏振纯度，但并非完美的线偏振光。
			
			最后，估算偏振度 \(P\)，使用公式 \(P = \frac{P_{\text{max}} - P_{\text{min}}}{P_{\text{max}} + P_{\text{min}}}\)。对于透射功率，计算得到 \(P_t = \frac{38.02 - 0.14}{38.02 + 0.14} \approx \frac{37.88}{38.16} \approx 0.993\)。接近 1 的偏振度表示高偏振纯度，但并非绝对完美。
		\end{itemize}
		综上所述，有限的消光比和非零的最小功率表明激光器发出的光不是完美的线偏振光。然而，高消光比和计算得到的偏振度表明激光器发出的光具有高度的线偏振性。因此，我们可以得出结论：该激光器发出的光不是完美的线偏振光，但具有高偏振纯度。
		
		此外，关于光轴得确定，借助了基于实验组内戴鹏辉同学思路构建的拟合代码进行分析。
		
		线偏振光经过半波片后的透射功率应满足公式 $P_t = E_0^2 \sin^2(2 \alpha - \theta)$，考虑采用拟合函数$ P_t = P_0 \sin^2(2\alpha - \theta) + C $，结果如\cref{fig:apl13exp1q3}所示。
		
		\begin{figure}
			\centering
			\includegraphics[width=0.8\linewidth]{APL1_3_Exp1_Q3.png}
			\caption{透射功率与理论计算拟合}
			\label{fig:apl13exp1q3}
		\end{figure}
		
		由于我们将 半波片刻度的 29° 定义为数据列中的 0°，该刻度即为半波片的光轴。
		
		\item \textbf{HWP2应该位于什么位置才可以确定地测量 PBS2 之前的光确实在 D 态(或垂直于 D 态)？}
		要确定 HWP2（半波片 2）应位于何种角度，才能确定地测量 PBS2 之前的光确实在 D 态（45° 偏振）或与 D 态垂直的态（135° 偏振），需要对实验装置的光学特性进行详细分析。
		
		首先，半波片的作用是将入射光的偏振方向旋转 \(2\theta\) 角度，其中 \(\theta\) 是入射光偏振方向与半波片快轴之间的夹角。经过 HWP2 后，光的偏振角度变为：
		
		\[
		\theta_{\text{出}} = \theta_{\text{入}} + 2\theta
		\]
		
		对于 D 态，入射光的偏振角度为 \(\theta_{\text{入}} = 45^\circ\)。为了使 PBS2 透过最大功率，经过 HWP2 后的偏振方向应与 PBS2 的透射轴（通常为水平偏振，0°）一致，即需要满足：
		
		\[
		\theta_{\text{出}} = n \times 180^\circ \quad (n \in \mathbb{Z})
		\]
		
		将 \(\theta_{\text{出}}\) 表达式代入，得到：
		
		\[
		45^\circ + 2\theta = n \times 180^\circ
		\]
		
		解出 \(\theta\)：
		
		\[
		2\theta = n \times 180^\circ - 45^\circ \\
		\theta = \frac{n \times 180^\circ - 45^\circ}{2}
		\]
		
		当 \(n = 1\) 时：
		
		\[
		\theta = \frac{180^\circ - 45^\circ}{2} = 67.5^\circ
		\]
		
		当 \(n = 2\) 时：
		
		\[
		\theta = \frac{360^\circ - 45^\circ}{2} = 157.5^\circ
		\]
		
		同理，可以得到更多满足条件的角度，如 247.5° 和 337.5°。
		
		从实验数据可以看出，在 HWP2 角度为 45°、135°、225° 和 315° 时，PBS2 的透射功率达到最大值。这意味着在这些角度下，经过 HWP2 后的光偏振方向与 PBS2 的透射轴对齐，验证了光确实处于 D 态。
		
		如果要确定光是否处于与 D 态垂直的态（即偏振角度为 135°），则需要使经过 HWP2 后的偏振方向与 PBS2 的透射轴正交。此时，偏振方向应为垂直偏振（90° 或 270°），因此需要满足：
		
		\[
		\theta_{\text{出}} = 90^\circ + n \times 180^\circ
		\]
		
		代入 \(\theta_{\text{出}} = 45^\circ + 2\theta\) 后，得到：
		
		\[
		45^\circ + 2\theta = 90^\circ + n \times 180^\circ
		\]
		
		解出 \(\theta\)：
		
		\[
		2\theta = 45^\circ + n \times 180^\circ \\
		\theta = \frac{45^\circ + n \times 180^\circ}{2}
		\]
		
		当 \(n = 0\) 时：
		
		\[
		\theta = \frac{45^\circ}{2} = 22.5^\circ
		\]
		
		当 \(n = 1\) 时：
		
		\[
		\theta = \frac{225^\circ}{2} = 112.5^\circ
		\]
		
		当 \(n = 2\) 时：
		
		\[
		\theta = \frac{405^\circ}{2} = 202.5^\circ
		\]
		
		因此，在 HWP2 角度为 22.5°、112.5°、202.5° 和 292.5° 时，PBS2 的透射功率达到最小值，表明光的偏振方向与 PBS2 的透射轴正交，光处于与 D 态垂直的态。
		
		综上所述，通过调整 HWP2 的角度，可以确定 PBS2 之前的光偏振态：
		\begin{itemize}
			\item 当 HWP2 设置为 45°、135°、225° 或 315° 时，PBS2 透射功率最大，确定光处于 D 态。
			\item 当 HWP2 设置为 22.5°、112.5°、202.5° 或 292.5° 时，PBS2 透射功率最小，确定光处于与 D 态垂直的态。
		\end{itemize}
	\end{itemize}
	
	\item 编写Python脚本计算误码率
	
	在实验过程中，我们编写了脚本（见\textbf{附件（Attachment）}部分）用于“手动”计算二进制误码率：
	
	首先，\texttt{hex\_to\_binary\_file}函数的作用是将一个指定的16进制文件转换为二进制文件。通过\texttt{open}函数，以二进制模式读取输入的16进制文件，并以二进制模式写入到输出文件。\texttt{hex\_f.read()}读取整个16进制文件的数据，然后使用\texttt{format(byte, '08b')}将每个字节转换为对应的8位二进制字符串。最后，二进制字符串被写入到指定的输出文件中。这个过程确保16进制文件的每一个字节都能被精确转换为二进制格式，便于后续的误码率计算。
	
	接下来是\texttt{calculate\_ber}函数，它用于计算两个二进制文件之间的误码率。函数首先以二进制模式分别打开发送端和接收端的二进制文件，读取所有的位数据。为了保证误码率的准确计算，程序首先检查两个文件的长度是否一致。如果两个文件的长度不同，函数将抛出一个错误提示，提醒无法进行误码率计算。接下来，程序通过遍历发送文件和接收文件的每一位数据，计算出其中不同的位数，也就是发生错误的位数。最后，误码率通过将错误的位数除以总的位数得到。
	
	在主程序中，首先定义了发送端和接收端的16进制文件路径，以及转换后的二进制文件路径。然后调用\texttt{hex\_to\_binary\_file}函数，将两个16进制文件转换为对应的二进制文件。最后，使用\texttt{calculate\_ber}函数计算这两个二进制文件之间的误码率，并将结果输出。整个流程确保了数据在传输过程中的准确性检测，误码率越低，表示传输质量越好。
\end{enumerate}

%% 总结
%\subsubsection{Conclusion}
%\begin{itemize}
%	\item 
%\end{itemize}


%---------------------------------------------------------------------
% 实验后思考题
\subsection{Reflections after Experiment}

%思考题1
\begin{question}
	是否可以通过直接衰减任意的光源（比如白炽灯）的强度到单光子级别来得到真正的单光子源？（真正的单光子源是指每次触发可以确定性的得到一个仅包含一个光子的光脉冲信号。）
\end{question}
答案是否定的，具体分析如下：

通过直接衰减任意光源（例如白炽灯）到单光子级别\textbf{无法}得到真正的单光子源。不同类型的光源具有不同的光子统计特性，这些统计特性决定了光子的发射模式。白炽灯等经典光源（热光源）的光子发射遵循\textbf{泊松统计}或\textbf{热统计}，而真正的单光子源具有更加严格的光子发射行为，遵循\textbf{定域性}和\textbf{确定性发射}的统计特性。对于经典光源，光子的发射遵循泊松分布，即使经过衰减，光子数的概率分布仍然呈泊松分布。这意味着即使在非常低的强度下，仍然有较大概率存在多个光子同时出现的现象。这种分布可以用以下公式描述：

\[
P(n) = \frac{\mu^n e^{-\mu}}{n!}
\]

其中，\(P(n)\) 是发射 \(n\) 个光子的概率，\(\mu\) 是平均光子数（光源的强度），\(n\) 是光子的数量。即便衰减到平均每个脉冲只有一个光子（\(\mu = 1\)），仍然有概率发射多个光子。例如，对于 \(\mu = 0.1\) 的情况下，仍然有一定概率（\(\sim 0.005\)）会发射两个或更多光子。这意味着使用经典光源和简单衰减的情况下，不可能得到确定性的单光子输出。而真正的单光子源具有\textbf{亚泊松统计}，确保每次脉冲只发射一个光子，其特点是发射0个光子的概率和发射超过1个光子的概率都为0。

在量子密钥分发（QKD）中，多光子态是一个安全隐患。如果存在多个光子被发射，窃听者（Eve）可以截取其中一个光子并测量其量子态，而不会影响其他光子的传输和测量结果。这种情况下，QKD协议的安全性会受到破坏，因为它依赖于量子测量的不可逆性。使用经典光源经过衰减到单光子水平时，虽然大多数脉冲可能包含0个或1个光子，但仍然有一定概率出现多光子事件。窃听者可以利用这种多光子态进行“光子分裂攻击”，从而破坏协议的安全性。因此，经典光源和简单衰减的方式无法确保密钥分发的安全性。

为了获得真正的单光子源，通常需要使用特殊的技术，例如量子点、色心等。这些技术能够精确地控制光子发射，确保每次脉冲只发射一个光子。量子点可以通过激发单个电子-空穴对，随后该对复合释放一个光子，达到单光子发射的目的。而色心光子源则利用特定的缺陷原子，通过激发态和基态之间的跃迁发射单个光子。这些技术都能够提供确定性的单光子发射，避免了多光子态的安全隐患。

即使通过光学衰减器调节经典光源的强度，发射的光子仍然遵循泊松分布，且发射多个光子的概率无法为零。假设经典光源的平均光子数为 \(\mu\)，使用衰减器将输出调节到单光子水平后，新的平均光子数为 \(\mu' = \mu \cdot T\)，其中 \(T\) 是衰减器的透过率。尽管可以将平均光子数接近1，发射多个光子的概率 \(P(n > 1)\) 始终存在。对于真正的单光子源，期望的发射条件为：\(P(n = 1) = 1, P(n > 1) = 0, P(n = 0) = 0\)。使用经过衰减的经典光源时，总会有一部分光子发射事件会产生多光子态，因此衰减经典光源无法提供确定性的单光子源。

% 思考题2
\begin{question}
	了解单光子探测器暗计数的原理，并设计实验装置测量探测器的暗计数率。
\end{question}
单光子探测器的暗计数是指在没有光子入射的情况下，探测器仍然产生误触发信号的现象。这种误触发通常来源于探测器本身的热噪声、材料缺陷、或环境背景辐射。暗计数率是指单位时间内探测器在无光条件下产生的错误探测事件的数量，通常以赫兹（Hz）为单位表示。暗计数的来源包括热电子噪声、环境辐射、以及探测器材料中的结构缺陷，这些因素都可能导致误探测，从而增加暗计数。特别是在雪崩光电二极管（APD）或光电倍增管（PMT）等单光子探测器中，由于其极高的灵敏度，暗计数会显著影响测量的准确性，尤其在量子密钥分发等实验中。

理论上，暗计数率（\( DCR \)）可以通过测量无光条件下探测器的触发次数来计算。假设在测量时间 \( T \) 内探测器产生了 \( N \) 次暗计数，则暗计数率的公式为：
\[
DCR = \frac{N}{T}
\]
为获得准确的暗计数率，通常需要在严格控制的无光条件下对探测器进行长时间测量，以避免由于短时间内偶然噪声引入误差。

测量单光子探测器暗计数率的实验装置设计需要满足完全无光的条件，以确保测量环境中没有任何光子入射。首先，实验装置包括单光子探测器，例如APD或PMT，这些探测器能够在极低光强下进行探测，并且对光子事件产生高度灵敏的响应。其次，需要一个暗室或光学屏蔽装置，用于屏蔽外界的环境光和背景噪声。这个暗室可以是一个全封闭的实验箱，确保内部没有任何光泄漏。探测器的输出端通过放大器与信号处理系统连接，信号处理系统的输出传送至计数器，以记录误触发事件的数量。此外，还需要一个时间控制器，用来精确控制测量的时间窗口，从而记录一段时间内的暗计数事件。实验方案设计如下：

\begin{ubox}{实验装置设计：测量单光子探测器的暗计数率}
	\highlight{实验目标}
	
	测量单光子探测器在无光子条件下的暗计数率，即单位时间内由于探测器热噪声、环境噪声等因素导致的误触发事件的频率。
	
	\highlight{实验装置组成}
	\begin{enumerate}
		\item 单光子探测器（APD或PMT）：单光子探测器是测量暗计数的核心组件，典型的选择包括雪崩光电二极管（APD）和光电倍增管（PMT）。这些探测器具有极高的灵敏度，能够探测到单个光子事件。为了保证实验的准确性，探测器需要在适当的工作条件下运行（如适当的偏压、温度等）。
		
		\item 暗室或光学屏蔽装置：暗室是用来隔绝任何外部光线的环境，确保探测器不会受到外部光源的干扰。暗室可以是一个完全封闭的光学箱或一个经过专门设计的实验室暗室，其内壁应覆盖不反光的材料以避免杂散光。暗室必须能够有效阻隔任何环境光或背景辐射的进入，确保探测器所处的环境完全无光。
		
		\item 信号处理系统：探测器输出的信号通常是电流或电压脉冲，这些信号需要通过信号处理系统进行放大和整形。信号处理系统的目的是将探测器产生的微弱信号放大为可测量的电信号，同时滤除不必要的噪声。这个系统通常包括前置放大器和信号整形器。
		
		\item 计数器：计数器用于记录探测器的输出脉冲数，即每次探测器误触发时产生的信号脉冲。计数器需要具有较高的时间分辨率和精确的计数能力，以确保能够准确统计每秒的暗计数事件。
		
		\item 时间控制器：时间控制器用于精确控制测量的时间窗口，并与计数器配合工作。测量暗计数率的时间窗口通常应较长（例如1分钟或更长），以确保统计结果的可靠性。时间控制器可以精确定义测量的起止时间，从而实现对实验过程的严格控制。
		
		\item 温度控制装置（可选）：由于探测器的暗计数率可能会受到温度的影响，实验装置中可以配备一个温度控制装置，以保持探测器在稳定的温度环境下工作。特别是对于雪崩光电二极管（APD），较低的温度可以有效降低热噪声和暗计数率。
	\end{enumerate}
	
	\highlight{实验步骤}
	\begin{enumerate}
		\item 装置准备  
		\begin{itemize}
			\item 将单光子探测器安装在暗室内部的适当位置，确保探测器的光学窗口不会暴露在任何外部光源下。
			\item 检查探测器的工作状态，调节探测器的工作偏压和温度（如果有温控装置），使探测器处于最佳工作条件。
			\item 连接探测器输出端到信号处理系统，确保探测器输出的电信号可以传送到计数器进行记录。
		\end{itemize}
		
		\item 系统初始化  
		\begin{itemize}
			\item 启动信号处理系统、计数器和时间控制器，确保所有设备正常工作。
			\item 确保暗室内没有任何光源进入，进行暗室封闭性检查，保证测量环境完全无光。
		\end{itemize}
		
		\item 开始测量  
		\begin{itemize}
			\item 设置时间控制器的测量窗口，例如1分钟或5分钟，确保能够获得足够长时间的统计数据。
			\item 启动计数器，开始记录探测器在此期间产生的误触发脉冲数，即暗计数事件。
			\item 记录计数器上的数据，即测量时间内探测器的误触发事件数 \(N\)。
		\end{itemize}
		
		\item 计算暗计数率：根据测量结果 \(N\) 和测量时间 \(T\)，使用公式计算暗计数率：
		\[
		DCR = \frac{N}{T}
		\]
		例如，如果在60秒的测量时间内记录到120次暗计数事件，则暗计数率为：
		\[
		DCR = \frac{120}{60} = 2 \, \text{Hz}
		\]
		这意味着探测器每秒在无光条件下产生2次误触发事件。  
		
		\item 多次重复测量：为了提高测量的准确性，可以多次重复测量，并取多次测量结果的平均值作为最终暗计数率。每次测量前都应确保探测器的工作状态和环境条件一致。  
	\end{enumerate}
\end{ubox}

% 思考题3
\begin{question}
	量子密钥分发实验中对基后的数据，为什么会有误码出现，如何降低误码？
\end{question}
\begin{enumerate}
	\item 在量子密钥分发（QKD）实验中，误码率（Quantum Bit Error Rate, QBER）是评估密钥分发质量和安全性的重要指标。尽管理论上QKD能够保证无条件的安全性，但在实际操作中，误码的出现是不可避免的，主要源自系统的物理限制和量子态测量中的不确定性。误码的来源包括探测器的暗计数、信道噪声、光子的多光子态问题、基选择不匹配以及探测器效率不一致等因素。
	
	暗计数是单光子探测器在没有光子信号时仍然产生的伪信号，会误导系统记录错误的偏振信息，从而导致误码。信道噪声是另一大误码来源，量子态在通过光纤或自由空间信道传输时，可能受到光纤散射、吸收以及大气湍流等环境因素的干扰，导致光子偏振态的错误测量。实验中使用的强衰减激光脉冲虽能模拟单光子源，但偶尔会产生多光子态，这不仅可能被窃听者截获，还会引发测量误差。基选择不匹配也是误码的常见原因，在BB84协议中，Alice和Bob随机选择基来制备和测量光子的偏振态，只有基匹配时才能正确测量光子，基的不匹配则会导致错误测量。探测器效率的不一致性也会引发误码，如果不同偏振态的光子无法被一致检测，就会导致错误的信号识别。
	
	\item 为了降低QKD系统中的误码率，首先可以通过冷却探测器或选用低暗计数率的探测器来减少暗计数。同时，改善量子信道的质量，如使用低损耗光纤、增加信道隔离措施以及通过实时偏振补偿来减少噪声引起的误码。增强光源的单光子性也是减少误码的有效手段，可以采用真正的单光子源如量子点或色心，或者优化激光脉冲的衰减，使得每个脉冲产生的光子数更接近单光子。此外，虽然BB84协议中基选择的不匹配是固有的，但可以通过改进协议（如SARG04）或使用多态量子编码来提高基匹配率。最后，误差纠正和隐私放大也是确保密钥安全的关键步骤，Alice和Bob可以通过经典通信纠正误码，并通过哈希函数进行隐私放大，消除可能被窃听的信息，确保最终密钥的安全性。
\end{enumerate}